\documentclass[11pt]{article}

\usepackage[T1]{fontenc}
\usepackage{lmodern}
\usepackage{amsmath,amssymb,amsthm,mathtools}
\usepackage{mathrsfs}
\usepackage{enumitem}
\usepackage[margin=1.12in]{geometry}
\usepackage{microtype}
\usepackage{xcolor}
\usepackage[colorlinks=true,linkcolor=blue!55!black,citecolor=blue!55!black,
  urlcolor=blue!55!black]{hyperref}
\emergencystretch=2em

\newtheorem{theorem}{Theorem}[section]
\newtheorem{proposition}[theorem]{Proposition}
\newtheorem{lemma}[theorem]{Lemma}
\newtheorem{corollary}[theorem]{Corollary}
\theoremstyle{definition}
\newtheorem{definition}[theorem]{Definition}
\newtheorem{example}[theorem]{Example}
\theoremstyle{remark}
\newtheorem{remark}[theorem]{Remark}

\newcommand{\C}{\mathbb C}
\newcommand{\R}{\mathbb R}
\newcommand{\N}{\mathbb N}
\newcommand{\Z}{\mathbb Z}
\newcommand{\cP}{\mathcal P}
\newcommand{\cT}{\mathcal T}
\newcommand{\cE}{\mathcal E}
\newcommand{\cS}{\mathcal S}
\newcommand{\cO}{\mathcal O}
\newcommand{\supp}{\operatorname{supp}}
\newcommand{\Span}{\operatorname{span}}
\newcommand{\Frac}{\operatorname{Frac}}
\newcommand{\trdeg}{\operatorname{trdeg}}
\newcommand{\Diff}{\partial}

\title{Projective-Tail Moduli of Signed Prime Dirichlet Fields:\\
Differential Matroids and Positive Source Pointing}
\author{Michael Haufschild \and GPT 5.6 Sol \and Fable 5 \and Grok 4.5}
\date{}

\begin{document}
\maketitle

\begin{abstract}
For an orientation \(\varepsilon:\mathcal P\to\{\pm1\}\) of the rational primes, put
\[
  P_\varepsilon(s)=\sum_{p\in\mathcal P}\varepsilon(p)p^{-s},
  \qquad \Re s>1.
\]
We classify differential-algebraic independence for every finite family of these functions and
their translates.  Let
\(\mathcal T=\C^{\mathcal P}/\C^{(\mathcal P)}\) be the quotient of prime
coefficient functions by finite support, and let \(\mathcal E_{\rm fin}\) be the differential
field generated by \(\C(s)\) and convergent Dirichlet series involving only finitely many
primes.  If \(\varepsilon_1,\ldots,\varepsilon_m\) are orientations and the finite shift set
has pairwise distinct real parts, then the complete translate family is differentially
algebraically independent over \(\mathcal E_{\rm fin}\) if and only if
\([\varepsilon_1],\ldots,[\varepsilon_m]\) are linearly independent in \(\mathcal T\).
Equivalently, the sign patterns recurring at infinitely many primes span \(\C^m\).

This yields an exact moduli theorem.  The differential fields generated by two complete
nonnegative-integer translate palettes agree precisely when their orientations are eventually
equal or eventually opposite.  Outside that projective tail relation the two palettes are jointly
differentially algebraically independent.  Hence the embedded field family is classified by the
Boolean tail quotient modulo finite changes and one global sign.  Every fibre of this quotient is
dense in the prime sign cube, so its quotient topology is indiscrete and every nonconstant
separated class invariant is nowhere continuous.

The Euler field fibre nevertheless has a unique positive source-level representative.  On every
line \(\Re s=\sigma>1\),
the vertical trace \(t\mapsto P_\varepsilon(\sigma+it)\) is positive definite exactly for the
all-positive orientation.  Thus differential-field type and positive Euler marking are transverse:
the former is a projective tail invariant, while the latter selects one representative and cannot
descend to the quotient.  The result concerns the first prime layer; it makes no zero-location
claim for the Riemann zeta function.
\end{abstract}

\noindent\textbf{Keywords.}
Dirichlet series; differential algebra; hypertranscendence; prime zeta function; tail quotient;
differential matroid; positive-definite function.

\medskip
\noindent\textbf{2020 Mathematics Subject Classification.}
11M41, 12H05, 13N15, 42A82.

\section{Introduction}

The classical hypertranscendence theorems for Gamma and zeta functions exclude finite algebraic
differential equations~\cite{Holder1887,Ostrowski1920}.  For Dirichlet series, Shapiro and Sparer
developed a systematic Jacobian method on the convolution algebra of arithmetic functions
\cite{ShapiroSparer1986}.  Pong subsequently proved a strong one-source theorem: if the prime
support of an arithmetic function is not finitely generated, then all of its logarithmic jets and
all translates with distinct real parts are algebraically independent over the finite-prime
coefficient algebra~\cite[Theorem~5.9]{Pong2016}.
We use differential-algebraic independence in the standard sense of Kolchin~\cite{Kolchin}.

Those results answer a freeness question for one source.  They do not classify simultaneous
families of different prime-supported sources, nor do they identify when two such sources generate
the same embedded differential field.  This paper addresses that multi-source problem for the
Boolean prime family
\begin{equation}\label{eq:intro-carrier}
  P_\varepsilon(s)=\sum_p\varepsilon(p)p^{-s},
  \qquad \varepsilon(p)\in\{\pm1\}.
\end{equation}

The answer is a positive constraint law.  Differential algebra sees exactly the projective tail
of the prime coefficients.  It forgets finitely many Euler coordinates and one simultaneous
global sign, but it forgets nothing else.

To state the result, let
\[
  \cT=\C^{\cP}/\C^{(\cP)}
\]
be the prime tail quotient.  Let \(\cE_{\rm fin}\) be the differential field of right-half-plane
meromorphic germs generated by \(\C(s)\) and by all convergent Dirichlet series whose supports
involve only finitely many primes.  For a finite family of orientations
\(\varepsilon=(\varepsilon_j)_{j\in J}\), write
\[
  \operatorname{Rec}(\varepsilon)
  =\left\{v\in\{\pm1\}^{J}:
       \#\{p:(\varepsilon_j(p))_{j\in J}=v\}=\infty\right\}.
\]

The main theorem has four parts.

\begin{theorem}[Prime-tail moduli theorem]\label{thm:intro-main}
Let \(J\) be finite and nonempty, let \(A\subset\C\) be finite and nonempty, and assume
that distinct elements of \(A\) have distinct real parts.
\begin{enumerate}[label=\textup{(\roman*)},leftmargin=2.4em]
  \item The family
  \[
    \bigl(P_{\varepsilon_j}(s+a)\bigr)_{j\in J,\,a\in A}
  \]
  is differentially algebraically independent over \(\cE_{\rm fin}\) if and only if
  \([\varepsilon_j]_{j\in J}\) is linearly independent in \(\cT\).  This is further
  equivalent to
  \[
    \Span_\C\operatorname{Rec}(\varepsilon)=\C^J.
  \]
  \item For orientations \(\varepsilon,\delta\), the differential fields generated by their
  complete nonnegative-integer translate palettes agree if and only if
  \(\varepsilon\) and \(\delta\) are eventually equal or eventually opposite.
  Outside this relation the two complete palettes are jointly differentially algebraically
  independent.
  \item The quotient of \(\{\pm1\}^{\cP}\) by eventual equality and eventual opposition has
  indiscrete quotient topology.  Every continuous class invariant into a \(T_0\) space is
  constant; every nonconstant class invariant into a \(T_1\) space is nowhere continuous.
  \item For every \(\sigma>1\), the function
  \(t\mapsto P_\varepsilon(\sigma+it)\) is positive definite on \(\R\) if and only if
  \(\varepsilon(p)=1\) for every prime \(p\).  In the projective tail class of the
  all-positive source, this is the unique positive-definite representative.
\end{enumerate}
\end{theorem}

Part (i) is the algebraic classifier.  Its analytic engine is classical: prime-basic derivations,
the Shapiro--Sparer Jacobian criterion, and uniqueness for convergent Dirichlet series.  The new
step is the converse multi-source identification of the entire differential matroid with linear
algebra in \(\cT\).  The finite sign range is crucial.  It turns a putative translate--jet
relation into an exponential-polynomial identity on an unbounded recurring-pattern fibre, where
distinct real parts separate the shift blocks.

Part (ii) promotes the matroid to an exact field moduli theorem.  Finite changes give finite-prime
corrections and global opposition gives simultaneous negation, so both preserve the generated
differential field.  The converse is much stronger: every pair outside those two alternatives is
jointly differentially algebraically independent.  No additional field identifications occur.

Parts (iii) and (iv) describe the resulting geometry.  The field quotient is canonical but
nonseparated: finite modifications make every fibre dense.  Positive definiteness does not descend
to that quotient, since changing one prime leaves the field fixed but destroys positivity.  It
instead gives a unique pointed lift of the Euler-positive class.  The differential field records
constraint type; positivity records the chosen source.

This distinction also clarifies the relation to the Riemann zeta function.  The all-positive
source is the prime zeta function, and on \(\Re s>1\)
\begin{equation}\label{eq:intro-plethysm}
  \log\zeta(s)=\sum_{r\geq1}\frac{1}{r}P_+(rs).
\end{equation}
The theorem classifies the first prime layer \(r=1\).  It neither controls the prime-power
plethysm in \eqref{eq:intro-plethysm} after continuation nor locates zeros.  Its RH-facing content
is therefore architectural: a zeta-sensitive constraint in this setting must retain the pointed
Euler representative and couple it to the prime-power and functional-equation data.

The paper is organized as follows.  Sections~\ref{sec:tail} and~\ref{sec:operator} establish the
tail and shift--clock laws.  Section~\ref{sec:transfer} carries them through the classical
arithmetic-function Jacobian method.  Section~\ref{sec:field-moduli} proves the exact field
quotient, Section~\ref{sec:topology} its nonseparated topology, and
Section~\ref{sec:pointing} its positive representative-level pointing.  The zeta boundary is
discussed in Section~\ref{sec:zeta}.

\section{Prime tails and recurring sign patterns}\label{sec:tail}

Let \(\cP\) be the set of rational primes and
\[
  \C^{(\cP)}
  =\{u:\cP\to\C:\supp u\text{ is finite}\}.
\]
The quotient
\begin{equation}\label{eq:tail-quotient}
  \cT=\C^{\cP}/\C^{(\cP)}
\end{equation}
is a complex vector space.  We write \([u]\) for the class of \(u\).

For a finite nonempty index set \(J\), let
\(\varepsilon_j:\cP\to\{\pm1\}\), \(j\in J\), and define the joint pattern
\[
  v_\varepsilon(p)=(\varepsilon_j(p))_{j\in J}\in\{\pm1\}^{J}.
\]
Its recurring patterns are
\begin{equation}\label{eq:recurring-patterns}
  \operatorname{Rec}(\varepsilon)
  =\{v\in\{\pm1\}^{J}:v_\varepsilon^{-1}(v)\text{ is infinite}\}.
\end{equation}

\begin{proposition}[Recurring-pattern criterion]\label{prop:recurring-pattern}
The following are equivalent:
\begin{enumerate}[label=\textup{(\roman*)},leftmargin=2.2em]
  \item the classes \([\varepsilon_j]\), \(j\in J\), are linearly independent in
  \(\cT\);
  \item \(\Span_\C\operatorname{Rec}(\varepsilon)=\C^J\).
\end{enumerate}
\end{proposition}

\begin{proof}
Suppose first that the recurring patterns do not span.  There is a nonzero linear form
\(\ell:\C^J\to\C\) vanishing on every recurring pattern.  Since the joint range
\(\{\pm1\}^{J}\) is finite, every nonrecurring fibre is finite.  Hence
\(p\mapsto\ell(v_\varepsilon(p))\) has finite support, giving a nontrivial relation among
the tail classes.

Conversely, suppose \(\ell(v_\varepsilon(p))\) has finite support.  On every infinite
recurring fibre it vanishes away from at most finitely many primes, and therefore it vanishes on
the corresponding pattern.  If the recurring patterns span \(\C^J\), then \(\ell=0\).
\end{proof}

The proposition makes the tail matroid concrete.  It depends not on densities of sign patterns
but only on which patterns occur infinitely often.

\begin{example}[Two orientations]\label{ex:two-orientations}
For two orientations \(\varepsilon,\delta\), their tail classes are dependent if and only if
\(\delta\) is eventually equal to \(\varepsilon\) or eventually equal to
\(-\varepsilon\).  Indeed, if
\(a\varepsilon+b\delta\) has finite support and \(a,b\neq0\), then
\(\delta(p)/\varepsilon(p)=-a/b\) for all sufficiently large primes.  The ratio takes only the
values \(\pm1\), so it is eventually constant.  Neither \(a\) nor \(b\) can vanish because an
orientation is nonzero at every prime.
\end{example}

\section{The shift--clock collapse}\label{sec:operator}

For \(a\in\C\), multiplication by \(p^{-a}=e^{-a\log p}\) preserves finite support and
therefore induces an automorphism
\[
  M_a[u]=[p\mapsto p^{-a}u(p)]
\]
of \(\cT\).  Multiplication by \(\log p\) induces an injective linear operator
\[
  L[u]=[p\mapsto(\log p)u(p)].
\]

For a finite shift set \(A\subset\C\) and \(m\geq0\), define the order-\(m\)
operator-tail rectangle
\begin{equation}\label{eq:operator-rectangle}
  \mathscr R_m(\varepsilon,A)
  =\bigl(M_aL^k[\varepsilon_j]\bigr)_{
       j\in J,\,a\in A,\,0\leq k\leq m}.
\end{equation}

The only analytic input required for the collapse is the following elementary separation lemma.

\begin{lemma}[Exponential-polynomial separation]\label{lem:exp-poly}
Let \(a_1,\ldots,a_r\in\C\) have pairwise distinct real parts, and let
\(Q_1,\ldots,Q_r\in\C[X]\).  If
\[
  \sum_{\nu=1}^r e^{-a_\nu x_n}Q_\nu(x_n)=0
\]
on an unbounded real sequence \(x_n\to+\infty\), then every \(Q_\nu\) is zero.
\end{lemma}

\begin{proof}
Order the indices so that
\(\Re a_1<\cdots<\Re a_r\).  After multiplication by \(e^{a_1x_n}\),
\[
  Q_1(x_n)+\sum_{\nu=2}^r e^{-(a_\nu-a_1)x_n}Q_\nu(x_n)=0.
\]
The sum tends to zero exponentially relative to every polynomial.  Hence
\(|Q_1(x_n)|\to0\), which is impossible unless \(Q_1=0\).  Remove the first term and
continue inductively.
\end{proof}

\begin{theorem}[Finite-range operator collapse]\label{thm:operator-collapse}
Let \(A\subset\C\) be finite and nonempty, with pairwise distinct real parts.  For a finite
family of prime orientations \((\varepsilon_j)_{j\in J}\), the following are equivalent:
\begin{enumerate}[label=\textup{(\roman*)},leftmargin=2.2em]
  \item \([\varepsilon_j]_{j\in J}\) is linearly independent in \(\cT\);
  \item for every \(m\geq0\), the rectangle
  \(\mathscr R_m(\varepsilon,A)\) is linearly independent in \(\cT\).
\end{enumerate}
\end{theorem}

\begin{proof}
The implication (ii)\(\Rightarrow\)(i) follows by choosing any \(a_0\in A\) and taking
the order-zero columns.  Multiplication by \(p^{-a_0}\) is invertible in the tail quotient.

Assume (i), fix \(m\), and suppose that a linear combination of the columns in
\eqref{eq:operator-rectangle} has finite support.  Thus, outside a finite set of primes,
\begin{equation}\label{eq:tail-relation}
  \sum_{j\in J}\sum_{a\in A}\sum_{k=0}^m
    c_{j,a,k}\varepsilon_j(p)p^{-a}(\log p)^k=0.
\end{equation}
Fix a recurring pattern \(v\in\operatorname{Rec}(\varepsilon)\).  Its prime fibre is
unbounded.  On that fibre, \eqref{eq:tail-relation} becomes
\[
  \sum_{a\in A}p^{-a}Q_{a,v}(\log p)=0,
\]
where
\[
  Q_{a,v}(X)=\sum_{k=0}^m
     \left(\sum_{j\in J}c_{j,a,k}v_j\right)X^k.
\]
Lemma~\ref{lem:exp-poly}, applied with \(x=\log p\), gives \(Q_{a,v}=0\) for every
\(a\).  Consequently, for each pair \((a,k)\), the linear form
\(v\mapsto\sum_jc_{j,a,k}v_j\) annihilates every recurring pattern.  By
Proposition~\ref{prop:recurring-pattern}, those patterns span \(\C^J\), so every
\(c_{j,a,k}\) is zero.
\end{proof}

The theorem is the rigidity step absent from a purely one-source statement.  The logarithmic
clock and every finite shift palette create no new matroid dependencies: for Boolean prime
coefficients, all operator-tail information collapses to the original tail classes.

\section{Differential-algebraic transfer}\label{sec:transfer}

We now transfer Theorem~\ref{thm:operator-collapse} to analytic Dirichlet series.  The machinery
in this section is classical; details are included to fix the coefficient field and the converse.

Let \(\mathscr A\) be the \(\C\)-algebra of arithmetic functions under Dirichlet convolution.
For \(b\in\mathscr A\), let \([\supp b]\) be the set of primes dividing some integer in
\(\supp b\), and set
\[
  \mathscr S
  =\{b\in\mathscr A:[\supp b]\text{ is finite}\}.
\]
For a prime \(q\), the basic derivation is
\begin{equation}\label{eq:prime-deriv}
  (\partial_qb)(n)=v_q(nq)b(nq).
\end{equation}
The logarithmic derivation is
\begin{equation}\label{eq:log-deriv}
  (\partial_Lb)(n)=(\log n)b(n).
\end{equation}
Under the formal Dirichlet transform, \(\partial_L\) corresponds to \(-d/ds\), while
multiplication of coefficients by \(n^{-a}\) corresponds to translation by \(a\)
\cite{Pong2016}.

\begin{lemma}[Finite evaluation packet]\label{lem:evaluation-packet}
Let \(I\) be finite, let \(X\) be a set, and let \(b_i:X\to\C\), \(i\in I\).  If no
nonzero linear combination of the \(b_i\) has finite support, then, after deleting any prescribed
finite subset \(E\subset X\), there are points \(x_r\in X\setminus E\), \(r\in I\), for
which
\[
  \det(b_i(x_r))_{r,i\in I}\neq0.
\]
\end{lemma}

\begin{proof}
The evaluation vectors \((b_i(x))_{i\in I}\), \(x\in X\setminus E\), span \(\C^I\).
Otherwise a nonzero linear functional annihilating their span would give a nontrivial combination
supported in \(E\).  Choose a basis from the spanning family.
\end{proof}

\begin{proposition}[Prime-supported Jacobian transfer]\label{prop:jacobian-transfer}
Let \((b_i)_{i\in I}\) be a finite family of prime coefficient functions whose classes are
linearly independent modulo finite support.  Extend each \(b_i\) to an arithmetic function
\(\widetilde b_i\) supported on the primes.  Then the \(\widetilde b_i\) are algebraically
independent over \(\Frac(\mathscr S)\).
\end{proposition}

\begin{proof}
Suppose there is a polynomial relation with coefficients in \(\Frac(\mathscr S)\).  Clear
denominators.  Only finitely many coefficients occur, and the union \(E\) of the prime divisors of
their supports is finite.  Lemma~\ref{lem:evaluation-packet} gives primes
\(q_r\notin E\) such that \(\det(b_i(q_r))\neq0\).

Every \(\partial_{q_r}\) annihilates every coefficient in the relation.  Since
\(\widetilde b_i\) is supported on the primes, its derivative is the scalar arithmetic function
\[
  \partial_{q_r}\widetilde b_i=b_i(q_r)\mathbf 1.
\]
Thus the Jacobian determinant of the \(\widetilde b_i\) relative to the selected prime-basic
derivations is a nonzero scalar.  The Shapiro--Sparer Jacobian criterion gives algebraic
independence over their common
kernel~\cite{ShapiroSparer1986}.  That kernel contains all the coefficients of the cleared
relation, a contradiction.
\end{proof}

The adjoining of the independent coordinate \(s\) is controlled by the following uniqueness
lemma.

\begin{lemma}[Polynomial-weighted Dirichlet uniqueness]
\label{lem:polynomial-dirichlet}
Let \(b_0,\ldots,b_d:\N_{\geq1}\to\C\) have Dirichlet series converging absolutely in a
common right half-plane.  If
\begin{equation}\label{eq:poly-dirichlet-zero}
  \sum_{r=0}^d s^r\sum_{n\geq1}b_r(n)n^{-s}=0
\end{equation}
throughout a right half-plane, then every \(b_r\) is zero.
\end{lemma}

\begin{proof}
Put \(Q_n(X)=\sum_{r=0}^db_r(n)X^r\).  If some \(Q_n\) is nonzero, let \(n_0\) be the
least such index.  Multiplying \eqref{eq:poly-dirichlet-zero} by \(n_0^x\), for large real
\(x\), gives
\[
  Q_{n_0}(x)+\sum_{n>n_0}Q_n(x)(n_0/n)^x=0.
\]
Choose a real \(\sigma\) in the common half-plane of absolute convergence.  The absolute value
of the tail is bounded by
\[
  x^dn_0^\sigma
  \left(\frac{n_0}{n_0+1}\right)^{x-\sigma}
  \sum_{r=0}^d\sum_{n>n_0}|b_r(n)|n^{-\sigma},
\]
which tends to zero.  Hence \(Q_{n_0}(x)\to0\), impossible for a nonzero polynomial.
\end{proof}

Let \(\mathscr M_\infty\) be the field of meromorphic germs on sufficiently far right
half-planes.  Define \(\cE_{\rm fin}\subset\mathscr M_\infty\) to be the differential field
generated by \(\C(s)\) and by the convergent Dirichlet transforms of elements of
\(\mathscr S\).  Each individual element of \(\cE_{\rm fin}\) involves only finitely many
primes, although there is no fixed finite prime set for the whole field.

\begin{theorem}[Prime-tail differential matroid]\label{thm:prime-tail-matroid}
Let \(J\) be finite and nonempty, let \(A\subset\C\) be finite and nonempty with pairwise
distinct real parts, and let \((\varepsilon_j)_{j\in J}\) be prime orientations.  The following
are equivalent:
\begin{enumerate}[label=\textup{(\roman*)},leftmargin=2.2em]
  \item the functions
  \(P_{\varepsilon_j}(s+a)\), \(j\in J\), \(a\in A\), are differentially
  algebraically independent over \(\cE_{\rm fin}\);
  \item \([\varepsilon_j]_{j\in J}\) is linearly independent in \(\cT\);
  \item \(\Span_\C\operatorname{Rec}(\varepsilon)=\C^J\).
\end{enumerate}
\end{theorem}

\begin{proof}
The equivalence (ii)\(\Leftrightarrow\)(iii) is
Proposition~\ref{prop:recurring-pattern}.  Assume (ii).  Differential algebraic independence is
equivalent to algebraic independence of every finite jet rectangle.  For a fixed order \(m\),
the prime coefficient columns of that rectangle are
\begin{equation}\label{eq:analytic-columns}
  p\longmapsto
  \varepsilon_j(p)p^{-a}(-\log p)^k,
  \qquad j\in J,\ a\in A,\ 0\leq k\leq m.
\end{equation}
Theorem~\ref{thm:operator-collapse} makes their tail classes linearly independent.  Proposition
\ref{prop:jacobian-transfer} makes the corresponding prime-supported arithmetic functions
algebraically independent over the finite-prime coefficient field.

The convergent Dirichlet transform is injective and sends \eqref{eq:analytic-columns} to the jets
of \(P_{\varepsilon_j}(s+a)\).  We spell out the coefficient-field step.  A putative polynomial
relation uses only finitely many elements of \(\cE_{\rm fin}\).  Express those elements as rational
functions of \(s\) and of finitely many convergent transforms of elements of \(\mathscr S\), and
clear all denominators; the resulting coefficient polynomial remains nonzero.  Derivatives and
products of the latter transforms are again transforms of
elements of \(\mathscr S\), on a common sufficiently far right half-plane.  After expansion, the
cleared relation therefore has the form
\[
  \sum_{r=0}^{d}s^r\mathcal D(c_r)(s)=0,
\]
where \(\mathcal D\) denotes the Dirichlet transform and each \(c_r\) is an
\(\mathscr S\)-coefficient polynomial in the arithmetic functions corresponding to
\eqref{eq:analytic-columns}.  Lemma~\ref{lem:polynomial-dirichlet} gives \(c_r=0\) for every
\(r\).  Proposition~\ref{prop:jacobian-transfer}, applied over \(\Frac(\mathscr S)\), then says
that every one of those coefficient polynomials was zero.  Thus the original relation was zero.
Every finite jet rectangle is algebraically independent, proving (i).

Conversely, if (ii) fails, choose constants \(c_j\), not all zero, such that
\(\sum_jc_j\varepsilon_j\) has finite support.  For any \(a_0\in A\),
\[
  \sum_jc_jP_{\varepsilon_j}(s+a_0)
\]
is a finite-prime Dirichlet polynomial and lies in \(\cE_{\rm fin}\).  This is a nontrivial
differential-linear relation, so (i) fails.
\end{proof}

\begin{remark}[Exact novelty boundary]\label{rem:novelty-boundary}
Pong's one-source theorem already implies translate--jet independence for any single orientation
\cite[Theorem~5.9]{Pong2016}.  Theorem~\ref{thm:prime-tail-matroid} is not presented as a new
Jacobian criterion or as a new proof of one-source hypertranscendence.  Its contribution is the
multi-source if-and-only-if classification and the complete tail matroid it exposes.
\end{remark}

\begin{corollary}[Universal hypertranscendence in the sign cube]\label{cor:all-ht}
For every prime orientation \(\varepsilon\), the function \(P_\varepsilon\) is
hypertranscendental over \(\cE_{\rm fin}\), and every finite set of its translates with distinct
real parts is differentially algebraically independent.
\end{corollary}

\begin{proof}
The single tail class \([\varepsilon]\) is nonzero because \(\varepsilon\) is nonzero at every
prime.  Apply Theorem~\ref{thm:prime-tail-matroid} with \(|J|=1\).
\end{proof}

\section{The exact differential-field quotient}\label{sec:field-moduli}

For an orientation \(\varepsilon\), define two embedded differential fields:
\begin{align*}
  L_\varepsilon
  &=\cE_{\rm fin}\langle P_\varepsilon(s)\rangle_\Diff,\\
  K_\varepsilon
  &=\cE_{\rm fin}\langle P_\varepsilon(s+n):n\in\N\rangle_\Diff.
\end{align*}
Here \(\langle\cdot\rangle_\Diff\) denotes differential-field generation inside
\(\mathscr M_\infty\).

\begin{definition}[Projective tail relation]\label{def:projective-tail}
For prime orientations \(\varepsilon,\delta\), write
\(\varepsilon\sim_{\rm pt}\delta\) if either
\[
  \{p:\varepsilon(p)\neq\delta(p)\}\text{ is finite}
\]
or
\[
  \{p:\varepsilon(p)\neq-\delta(p)\}\text{ is finite}.
\]
\end{definition}

This is the equivalence relation generated by finite coordinate changes and one global
complement.

\begin{theorem}[Exact projective-tail field classification]
\label{thm:field-classification}
For any two prime orientations \(\varepsilon,\delta\), the following are equivalent:
\begin{enumerate}[label=\textup{(\roman*)},leftmargin=2.2em]
  \item \(\varepsilon\sim_{\rm pt}\delta\);
  \item \(L_\varepsilon=L_\delta\);
  \item \(K_\varepsilon=K_\delta\).
\end{enumerate}
If these conditions fail, then the two complete palettes
\[
  \bigl(P_\varepsilon(s+n)\bigr)_{n\geq0},
  \qquad
  \bigl(P_\delta(s+n)\bigr)_{n\geq0}
\]
are jointly differentially algebraically independent over \(\cE_{\rm fin}\).
\end{theorem}

\begin{proof}
If \(\varepsilon\) and \(\delta\) are eventually equal, then
\(P_\varepsilon-P_\delta\) is a finite-prime Dirichlet polynomial.  If they are eventually
opposite, then \(P_\varepsilon+P_\delta\) is such a polynomial.  The same statement holds after
every nonnegative-integer translation and every derivative.  Thus (i) implies both (ii) and (iii).

Suppose (i) fails.  By Example~\ref{ex:two-orientations}, the two tail classes
\([\varepsilon]\) and \([\delta]\) are linearly independent.  Theorem
\ref{thm:prime-tail-matroid} shows that the two unshifted carriers are jointly differentially
algebraically independent, and hence their generated differential fields cannot agree.  This
proves the contrapositive of (ii)\(\Rightarrow\)(i).

Every finite subset of the two complete natural-shift palettes is covered by
Theorem~\ref{thm:prime-tail-matroid}, because distinct natural shifts have distinct real parts.
Thus the union of the two complete palettes is differentially algebraically independent.  If
\(K_\varepsilon=K_\delta\), one palette element would be a differential rational function of
finitely many elements of the other palette, contradicting that independence.  This proves the
contrapositive of (iii)\(\Rightarrow\)(i) and the final assertion.
\end{proof}

\begin{corollary}[The field moduli set]\label{cor:field-moduli}
Let \(\Omega=\{\pm1\}^{\cP}\).  The assignment
\[
  \varepsilon\longmapsto K_\varepsilon
\]
induces a bijection from the projective tail quotient \(\Omega/{\sim_{\rm pt}}\) onto the
family of embedded differential fields \(\{K_\varepsilon:\varepsilon\in\Omega\}\).
The same statement holds with \(L_\varepsilon\) in place of \(K_\varepsilon\).
\end{corollary}

The word ``embedded'' matters.  The theorem classifies equality inside the common analytic germ
field, not merely abstract isomorphism types.  Its strength is that every failure of projective
tail equivalence produces joint differential independence, rather than only unequal fields.

\section{Nonseparated projective-tail topology}\label{sec:topology}

Equip \(\Omega=\{\pm1\}^{\cP}\) with the product topology and each sign coordinate with the
discrete topology.  After enumerating the primes, this is the Cantor sign cube.  Let
\[
  q:\Omega\longrightarrow\Omega/{\sim_{\rm pt}}
\]
be the quotient map.

\begin{theorem}[Dense field fibres and indiscrete quotient]\label{thm:indiscrete}
Every fibre of \(q\) is dense in \(\Omega\).  Consequently the quotient topology on
\(\Omega/{\sim_{\rm pt}}\) is indiscrete.
\end{theorem}

\begin{proof}
A basic open cylinder prescribes signs at only finitely many primes.  Given an orientation
\(\varepsilon\) and such a cylinder, overwrite those finitely many coordinates of
\(\varepsilon\) with the prescribed signs.  The resulting orientation lies in the cylinder and
is eventually equal to \(\varepsilon\).  Thus the ordinary finite-change orbit of
\(\varepsilon\), and hence its projective tail fibre, meets every nonempty basic open set.

If \(U\) is a nonempty open subset of the quotient, then \(q^{-1}(U)\) is a nonempty open
saturated subset of \(\Omega\).  Every fibre is dense and therefore meets this open set.
Saturation then forces \(q^{-1}(U)\) to contain every fibre, so \(q^{-1}(U)=\Omega\) and \(U\)
is the whole quotient.
\end{proof}

\begin{corollary}[Separated invariants collapse]\label{cor:continuous-collapse}
Let \(Y\) be a topological space and \(F:\Omega\to Y\) be constant on projective tail
classes.
\begin{enumerate}[label=\textup{(\roman*)},leftmargin=2.2em]
  \item If \(Y\) is \(T_0\) and \(F\) is continuous, then \(F\) is constant.
  \item If \(Y\) is \(T_1\) and \(F\) is continuous at even one orientation, then
  \(F\) is constant.
  \item If \(Y\) is \(T_1\) and \(F\) is nonconstant, then \(F\) is nowhere
  continuous.
\end{enumerate}
\end{corollary}

\begin{proof}
For (i), \(F\) factors through a continuous map from the indiscrete quotient to \(Y\).  A
continuous map from an indiscrete space to a \(T_0\) space is constant.

For (ii), suppose \(F\) is continuous at \(\varepsilon\) and that
\(F(\delta)\neq F(\varepsilon)\).  Since \(Y\) is \(T_1\), the complement of
\(\{F(\delta)\}\) is an open neighborhood of \(F(\varepsilon)\).  Continuity at
\(\varepsilon\) gives a neighborhood \(U\) of \(\varepsilon\) whose image avoids
\(F(\delta)\).  The projective tail fibre of \(\delta\) is dense, so it meets \(U\); class
invariance then gives a point of \(U\) with value \(F(\delta)\), a contradiction.  Part (iii)
is the contrapositive pointwise form of (ii).
\end{proof}

By Corollary~\ref{cor:field-moduli}, these are statements about the exact field fibres.  The
differential-field classifier is algebraically complete but admits no nonconstant continuous
separated coordinate on the coefficient cube.  This is not a defect in the classification; it is
the topology forced by finite Euler gauge invariance.

\section{Positive-definite source pointing}\label{sec:pointing}

The field quotient forgets the exact representative.  The convergence half-plane nevertheless
contains a canonical representative-level marker.

Recall that a function \(f:\R\to\C\) is positive definite if, for every finite choice of
\(t_1,\ldots,t_N\in\R\) and \(c_1,\ldots,c_N\in\C\),
\[
  \sum_{i,j=1}^N\overline{c_i}c_j f(t_j-t_i)\geq0.
\]

For \(\sigma>1\), define the vertical trace
\begin{equation}\label{eq:vertical-trace}
  T_{\varepsilon,\sigma}(t)
  =P_\varepsilon(\sigma+it)
  =\sum_p\varepsilon(p)p^{-\sigma}e^{-it\log p}.
\end{equation}
The series is the Fourier--Stieltjes transform of the finite signed atomic measure
\begin{equation}\label{eq:signed-measure}
  \mu_{\varepsilon,\sigma}
  =\sum_p\varepsilon(p)p^{-\sigma}\delta_{-\log p}.
\end{equation}

\begin{theorem}[Unique positive-definite orientation]\label{thm:positive-pointing}
For every \(\sigma>1\), the following are equivalent:
\begin{enumerate}[label=\textup{(\roman*)},leftmargin=2.2em]
  \item \(T_{\varepsilon,\sigma}\) is positive definite on \(\R\);
  \item \(\varepsilon(p)=1\) for every prime \(p\).
\end{enumerate}
\end{theorem}

\begin{proof}
If all signs are positive, \eqref{eq:signed-measure} is a finite positive measure, so its Fourier
transform is positive definite by Bochner's theorem~\cite{Bochner1933,RudinFourier}.

Conversely, if \(T_{\varepsilon,\sigma}\) is positive definite, Bochner's theorem represents it
as the Fourier transform of a finite positive measure.  Uniqueness of Fourier--Stieltjes transforms
of finite measures identifies that measure with \(\mu_{\varepsilon,\sigma}\).  The atoms
\(-\log p\) are distinct, so every mass \(\varepsilon(p)p^{-\sigma}\) must be nonnegative.
Since \(p^{-\sigma}>0\), every sign is positive.
\end{proof}

\begin{corollary}[Unique pointed lift of the Euler class]\label{cor:pointed-lift}
Let \(\varepsilon_+\) be the all-positive orientation.  Among all orientations
\(\varepsilon\) satisfying \(K_\varepsilon=K_{\varepsilon_+}\), exactly one has a
positive-definite vertical trace on any, equivalently every, line \(\Re s=\sigma>1\): namely
\(\varepsilon=\varepsilon_+\).
\end{corollary}

\begin{proof}
Theorem~\ref{thm:field-classification} says that the field fibre consists of orientations that are
eventually all positive or eventually all negative.  Theorem~\ref{thm:positive-pointing} selects
exactly the orientation that is positive at every prime.
\end{proof}

Thus pointing is not a function on field moduli.  A single sign flip leaves the differential field
unchanged but changes the positive-definite truth value.  The all-negative orientation generates
the same differential field as the all-positive orientation and lies at the opposite extreme of
the positivity cone.  Differential-field type and source orientation are therefore provably
transverse data.

\section{The zeta boundary}\label{sec:zeta}

Let \(\varepsilon_+(p)=1\).  Then
\[
  P_+(s)=\sum_pp^{-s}
\]
is the prime zeta function~\cite{Froberg1968}.  Corollary~\ref{cor:all-ht} makes it
hypertranscendental over \(\cE_{\rm fin}\), and
Theorem~\ref{thm:positive-pointing} identifies it as the unique positive-definite vertical source
in the entire sign cube.

The relation to the Riemann zeta function is obtained by expanding the Euler product in its
absolute-convergence half-plane:
\begin{equation}\label{eq:euler-plethysm}
  \log\zeta(s)
  =\sum_p\sum_{r\geq1}\frac{p^{-rs}}{r}
  =\sum_{r\geq1}\frac{P_+(rs)}{r},
  \qquad \Re s>1.
\end{equation}
The theorem therefore determines the differential matroid and the embedded-field fibres of the
first prime layer under Boolean reorientation, but not the compatibility among all prime-power
layers.

Accordingly, a zeta-sensitive extension of this first-layer model has to supply three kinds of
additional data.

\begin{enumerate}[label=\textup{(\arabic*)},leftmargin=2.3em]
  \item It must be \emph{source marked} if it is to recover the positive Euler orientation.  The
  embedded differential field alone cannot do so, because finitely changed and globally negated
  orientations generate the same field.  The positive-definite trace supplies such a marking on
  \(\Re s>1\).
  \item It must be \emph{prime-power complete} if it is to reconstruct \(\log\zeta\) through
  \eqref{eq:euler-plethysm}.  The first layer \(P_+(s)\) does not by itself encode the coupling
  over \(r\geq2\).
  \item It must be \emph{continuation sensitive} if it is to address zero location.  The tail
  theorem is proved in the absolute-convergence region; a zero-location statement requires data
  surviving analytic continuation and coupling to the functional equation.
\end{enumerate}

These are positive structural requirements, not an RH criterion.  The present theorem neither
implies the Riemann hypothesis nor supplies a mechanism that distinguishes the critical line.
It identifies which information is present in the hypertranscendental first-prime layer and which
information must be added before zero geometry can enter.

\section{Conclusion}

The signed prime Dirichlet family has an exact differential-algebraic moduli law.  For every finite
multi-source translate palette with distinct shift real parts,
\[
  \text{differential algebraic independence}
  \quad\Longleftrightarrow\quad
  \text{linear independence in }
  \C^{\cP}/\C^{(\cP)}.
\]
Because sign patterns have finite joint range, the complete shift--clock operator matroid collapses
to the original tail matroid.  The classical arithmetic-function Jacobian method then transfers
this coefficient law to analytic Dirichlet series.

The resulting field quotient is exact: finite prime changes and one global sign are all the
identifications.  Its topology is necessarily nonseparated because every field fibre is dense.
Positive definiteness supplies the missing representative-level datum and uniquely selects the
all-positive prime source.  Thus the central geometry is a pointed nonseparated moduli space:
projective tails classify embedded differential fields, while positivity chooses a source above
one field class.

For zeta, the theorem isolates the first-prime constraint geometry and shows exactly where an
extension must add positive source marking, prime-power coupling, and continuation data before
zero geometry can enter.

\begin{thebibliography}{99}

\bibitem{Bochner1933}
S.~Bochner,
Monotone Funktionen, Stieltjessche Integrale und harmonische Analyse,
\emph{Math. Ann.} \textbf{108} (1933), 378--410.

\bibitem{Froberg1968}
C.-E. Fr\"oberg,
On the prime zeta function,
\emph{BIT} \textbf{8} (1968), 187--202.

\bibitem{Holder1887}
O.~H\"older,
\"Uber die Eigenschaft der Gammafunction, keiner algebraischen Differentialgleichung zu
gen\"ugen,
\emph{Math. Ann.} \textbf{28} (1887), 1--13.

\bibitem{Kolchin}
E.~R. Kolchin,
\emph{Differential Algebra and Algebraic Groups},
Pure and Applied Mathematics, vol.~54, Academic Press, New York, 1973.

\bibitem{Ostrowski1920}
A.~Ostrowski,
\"Uber Dirichletsche Reihen und algebraische Differentialgleichungen,
\emph{Math. Z.} \textbf{8} (1920), 241--298.

\bibitem{Pong2016}
W.~Y. Pong,
Applications of differential algebra to algebraic independence of arithmetic functions,
\emph{Acta Arith.} \textbf{172} (2016), no.~2, 149--173.

\bibitem{RudinFourier}
W.~Rudin,
\emph{Fourier Analysis on Groups},
Interscience Tracts in Pure and Applied Mathematics, no.~12, Wiley, New York, 1962.

\bibitem{ShapiroSparer1986}
H.~N. Shapiro and G.~H. Sparer,
On algebraic independence of Dirichlet series,
\emph{Comm. Pure Appl. Math.} \textbf{39} (1986), no.~6, 695--745.

\end{thebibliography}

\end{document}
